% chapters/02/rangkuman.tex
\bagian{Rangkuman}

\begin{rangkumanbox}
\begin{itemize}
  \item Permainan dipandang sebagai sistem buatan yang memuat empat unsur:
        pemain, aturan, konflik atau tantangan, dan keluaran terukur.
  \item Program permainan adalah sistem waktu nyata: ia berjalan terus dan
        harus menanggapi masukan pemain dengan segera, sehingga pekerjaannya
        disusun sebagai siklus masukan--pembaruan--penayangan yang berulang.
  \item Program permainan tersusun dari subsistem masukan, pengendali waktu,
        simulasi, aturan dan tabrakan, penyajian, audio, serta data dan
        penyimpanan.
  \item Waktu tidak boleh diukur dengan jumlah bingkai. Kecepatan dinyatakan
        dalam satuan per detik, lalu dikalikan selang waktu
        ($\text{perpindahan} = \text{kecepatan} \times \Delta t$) agar perilaku
        permainan sama pada setiap komputer.
  \item Langkah berubah mengikuti lama bingkai dan paling mudah diterapkan;
        langkah tetap memakai panjang langkah yang sama dengan bantuan
        penampung waktu, sehingga hasilnya dapat diulang dan lebih stabil.
  \item Koordinat layar ber titik asal di kiri-atas dengan sumbu tegak menurun,
        berbeda dari koordinat kartesius. Selain itu, skala dan satuan harus
        ditetapkan secara sengaja.
  \item Keadaan gerak dinyatakan oleh posisi, kecepatan, dan percepatan.
        Tabrakan diperiksa dengan persegi pembatas sejajar sumbu, yang murah
        dihitung tetapi kurang teliti.
  \item Komputer mensimulasikan gerak secara terputus melalui kaidah Euler.
        Urutan pembaruannya penting, dan galatnya berbanding lurus dengan
        panjang langkah.
  \item Pengendalian NPC terbagi menjadi lapisan keputusan, yang dijawab
        dengan mesin keadaan, dan lapisan navigasi, yang dijawab dengan
        pencarian jalur pada graf.
  \item Pustaka menunggu dipanggil, sedangkan mesin memanggil kode Anda.
        Mesin mempercepat pengembangan tetapi menyembunyikan sebab akibat,
        sehingga pemahaman tentang mekanisme dasar tetap diperlukan.
\end{itemize}
\end{rangkumanbox}
